\section{Case Studies}

In this section, we discuss some case studies that we performed on
Haskell's libaries. We collected 25 functions from \hskin{Data.List}
and 10 frunctions from \hskin{Data.Map} library to test (\abha{what?})

\subsection{Data.List}
\begin{figure}
    \begin{subfigure}[t]{1\linewidth}
        \begin{smallhskcode}
appendTest = test2 (++)
    [ex| b*{a} |]  [ex| d*{c} |] .== [ex| b*{a} d*{c} |]
        \end{smallhskcode}
        \caption{\hskin{append}}
        \label{fig:list-append-ex}
    \end{subfigure}

\begin{subfigure}[t]{1\linewidth}
        \begin{smallhskcode}
filterTest = test2 (\x -> filter (/= x))
    [ex| x |] [ex| (x |{b} (y /= x))*{n} |]
                    .== [ex| ([] |{b} y)*{n} |]
        \end{smallhskcode}
        \caption{\hskin{filter}}
        \label{fig:list-filter-ex}
    \end{subfigure}
    \caption{\hskin{Data.List} example tests in \ourDSL}
    \label{fig:list-ex}
\end{figure}
Figure~\ref{fig:list-append-ex} shows the test written to generate 
input-output examples for the \hskin{append} function from \hskin{Data.List}. 
To clarify the process, \hskin{test2} serves as a higher-order 
function accepting three arguments: (1) the function under test, 
(2) a core expression representing the input arguments, 
and (3) a core expression representing the output. 

In the append example, the first argument to test2 is \hskin{(++)}, which is 
the list concatenation function of type \hskin{[Int] -> [Int] -> [Int]}. 
The second argument, \hskin;[ex| b*{a} |];, is the core expression for the 
first input to \hskin{(++)}, generating a random list of integers of length a. 
Similarly, the third argument \hskin;[ex| d*{c} |]; generates a random list 
of integers of length c. Lastly, the core expression \hskin;[ex| b*{a} d*{c} |]; 
stands for the concatenated output of these two lists.

\mypara{Higher-order functions}

\subsection{Data.Map}
\begin{figure}
    \begin{hskcode}
        lookupTest :: TestResult
        lookupTest = test2
            ((\x -> lookup x . fromList))
            [ex| k |]
            [ex| <((. /= k), .)>* <(k,v -> v)> <((. /= k), .)>* |] 
            .==
            [ex| <Just v> |]
    \end{hskcode}
    \caption{\hskin{lookup} test in \tool}
    \label{fig:lookup-test}
\end{figure}